% !TEX root = ../main.tex
\chapter{技术实施方案}

\section{技术选型原则}

本项目的技术选型遵循以下原则：
\begin{enumerate}
    \item \textbf{成熟稳定}：优先选择经过工业验证的技术，降低技术风险；
    \item \textbf{云原生友好}：支持容器化和无服务器部署，便于弹性伸缩；
    \item \textbf{性能优先}：在满足功能需求的前提下，优化运行效率。
\end{enumerate}

\section{技术栈概览}

基于上述原则，项目技术栈如表~\ref{tab:tech_stack_simple}所示。

\begin{table}[htbp]
    \centering
    \caption{项目技术栈}
    \label{tab:tech_stack_simple}
    \begin{tabular}{l l}
        \toprule
        \textbf{层次} & \textbf{技术选型}                                \\
        \midrule
        前端          & Flutter Web + fl\_chart + Firebase Auth      \\
        后端          & Python 3.11 + Flask 2.x                      \\
        ML/优化       & LightGBM + Gurobi + SHAP                     \\
        云服务         & Google App Engine + Firebase + Cloud Storage \\
        \bottomrule
    \end{tabular}
\end{table}

\section{系统架构}

系统采用\textbf{分层架构}设计，如图~\ref{fig:system_arch_simple}所示，自顶向下分为表现层、应用层、领域层和基础设施层。

\begin{figure}[H]
    \centering
    \begin{tikzpicture}[
            node distance=0.5cm,
            layer/.style={rectangle, draw, minimum width=12cm, minimum height=1.2cm, align=center, font=\small}
        ]
        % 分层
        \node[layer, fill=blue!15] (presentation) {表现层：Flutter Web + REST API + Firebase Auth};
        \node[layer, fill=green!15, below=of presentation] (application) {应用层：优化调度服务 + 负载预测服务 + 数据分析服务};
        \node[layer, fill=orange!15, below=of application] (domain) {领域层：MIP模型 + ML模型 + 特征工程};
        \node[layer, fill=purple!15, below=of domain] (infrastructure) {基础设施层：Firestore + Cloud Storage + 外部API};
    \end{tikzpicture}
    \caption{系统分层架构图}
    \label{fig:system_arch_simple}
\end{figure}

\section{云端部署架构}

系统部署在Google Cloud Platform上，采用全托管的PaaS服务，实现自动扩缩容与高可用。主要部署组件如表~\ref{tab:deploy_components}所示。

\begin{table}[H]
    \centering
    \caption{云端部署组件}
    \label{tab:deploy_components}
    \begin{tabular}{l l l}
        \toprule
        \textbf{组件} & \textbf{服务}             & \textbf{职责}    \\
        \midrule
        后端API服务     & Google App Engine       & Flask应用托管、自动扩缩 \\
        前端静态资源      & Firebase Hosting        & 全球CDN分发、HTTPS  \\
        用户认证        & Firebase Authentication & 身份验证、会话管理      \\
        数据存储        & Cloud Firestore         & 用户数据、操作历史      \\
        文件存储        & Cloud Storage           & CSV数据、ML模型文件   \\
        定时任务        & Cloud Scheduler         & 数据采集、模型更新触发    \\
        \bottomrule
    \end{tabular}
\end{table}

\section{部署流程}

项目采用CI/CD自动化部署流程，如图~\ref{fig:deploy_flow_simple}所示。

\begin{figure}[htbp]
    \centering
    \begin{tikzpicture}[
            node distance=1.5cm,
            box/.style={rectangle, draw, rounded corners, minimum width=2cm, minimum height=0.8cm, align=center, font=\footnotesize},
            arrow/.style={->, >=stealth, thick}
        ]
        \node[box, fill=gray!20] (code) {代码提交\\(Git Push)};
        \node[box, fill=blue!20, right=of code] (build) {构建镜像\\(Cloud Build)};
        \node[box, fill=green!20, right=of build] (deploy) {部署实例\\(GAE Deploy)};
        \node[box, fill=orange!20, right=of deploy] (route) {流量切换\\(Traffic Split)};

        \draw[arrow] (code) -- (build);
        \draw[arrow] (build) -- (deploy);
        \draw[arrow] (deploy) -- (route);
    \end{tikzpicture}
    \caption{CI/CD 部署流程}
    \label{fig:deploy_flow_simple}
\end{figure}

部署命令示例：
\begin{itemize}
    \item 后端部署：\texttt{gcloud app deploy back/app.yaml}
    \item 前端部署：\texttt{firebase deploy --only hosting}
\end{itemize}
